\section{GlyphChain security, recovery and bounded soak qualification}
\label{sec:glyphchain-security-recovery-soak}

This stage extends the production-security foundation with a repeatable,
content-free soak exercise.  It does not replace an independent application or
consensus review, and it does not convert a local machine into a production
network.  Its function is to prove that the durable control paths can be driven
repeatedly under bounded restart, interruption, redelivery, partition and
multi-network conditions without losing their safety invariants.

\subsection{Relationship to account recovery and production security}

Tessaris already has a monotonic application-account descriptor registry, a
fail-closed HTTPS client, account-root-signed descriptors, recovery epochs,
authorised-device checks and body-bound desktop sessions.  Production account
recovery nevertheless remains gated on deployment of the registry under
independent operational control, configuration of its production HTTPS URL,
connection of every application-session issuer to current-epoch enforcement,
and a clean-device lost-device exercise against that live service.

Validator-finality credentials remain deployment-managed secrets.  They are
not placed in an application-account recovery kit, copied during clean-device
recovery or granted to a recovered user device.  Application recovery restores
application authority only; it cannot confer consensus authority.

The new soak exercise addresses a different part of the release gate: repeated
local operation of the network control plane.  The two assurance statements
must not be merged.  Passing a network soak does not prove live account
revocation, and passing local account-recovery tests does not prove network
availability.

\subsection{Bounded multi-network topology}

The runner creates three isolated logical networks.  Each network receives a
distinct chain identifier, relay store and partition-recovery store.  The
qualification run performs 256 cycles per network, for 768 total message
cycles, and reconstructs the durable relay and recovery objects every sixteen
cycles.

The configured ranges are deliberately bounded:

\begin{itemize}
  \item two to eight logical networks;
  \item eight to 5,000 cycles per network;
  \item a restart interval between one and the selected cycle count;
  \item at most thirty-two queued relay items;
  \item at most sixteen queued items for one target;
  \item a 4,096-byte envelope cap; and
  \item a 120-second maximum relay time to live.
\end{itemize}

Invalid values fail before any exercise state is created.  The state root must
be empty, and a symbolic-link root is refused.

\subsection{Intermittent delivery and restart semantics}

Each cycle enqueues a bounded \texttt{BLOCK\_ANNOUNCE} envelope for an
intermittent recipient.  An immediate repeated put must resolve to the same
content-derived relay identifier and be reported as an idempotent duplicate.
The recipient then pulls the item twice before acknowledging it.  The second
pull must show a delivery-attempt count of two, and the acknowledgement must
remove exactly one item.

At each restart boundary, the in-memory relay object is discarded and rebuilt
from its durable state.  The test fails if the outstanding item is missing,
duplicated or addressed to a different recipient after reconstruction.  Every
network must finish with an empty queue.  This exercises at-least-once transport
with recipient acknowledgement while preventing duplicate queue insertion; it
does not claim that every downstream consensus handler is itself exactly once.

Only protocol metadata is admitted to the exercise.  The generated evidence
checks that document bodies, signer addresses, Legal matters and attachments do
not occur in the relay payloads.

\subsection{Repeated partition recovery}

Every message cycle drives the node through a bounded interruption-recovery
sequence:

\begin{equation}
  \texttt{HEALTHY}
  \longrightarrow \texttt{BLOCK\_CATCHUP}
  \longrightarrow \texttt{HEALTHY}.
\end{equation}

The transition journal is hash chained and persisted.  Halfway through each
network run, the exercise additionally introduces a conflicting trusted
finality condition.  This forces
\texttt{SAFETY\_HALT}; ordinary health observations cannot clear it.  The
exercise supplies the exact halt identifier, named operator, retained reason
and exact administrative confirmation, after which the node returns through
\texttt{REJOIN\_ASSESSING} and \texttt{BLOCK\_CATCHUP}.  Each network must end
healthy with a valid journal.

For an event body $E_i$, the journal link is

\begin{equation}
  H_i = \operatorname{SHA256}\!\left(
    \operatorname{canon}(E_i,H_{i-1})
  \right).
\end{equation}

Changing a recorded phase, height, reason or predecessor therefore invalidates
the chain.

\subsection{Cross-network isolation}

After the three soak loops, an envelope carrying the chain identifier of the
third network is submitted to the first network's relay.  Acceptance would be a
qualification failure.  The expected result is the explicit
\texttt{nat\_relay\_envelope\_chain\_mismatch} rejection.  This proves the
local route is chain-bound under the tested implementation; it does not prove
isolation between separately administered public infrastructures.

\subsection{Content-free evidence and tamper detection}

The exercise writes an owner-only JSON report with exclusive-create semantics.
It cannot overwrite an earlier report, and a symbolic-link destination is
refused.  The report contains configuration, aggregate counters, one-way
network commitments, final queue depths, partition-journal results, the
explicit assurance boundary and remaining external gates.  It contains no
document, signer, delivery attachment or Legal-matter data.

The report digest is

\begin{equation}
  H_R = \operatorname{SHA256}\!\left(
    \operatorname{canon}(R \setminus \{\texttt{report\_sha256}\})
  \right).
\end{equation}

The verifier recomputes $H_R$ using constant-time comparison.  Modifying a
counter, check, configuration value, claim boundary or remaining gate therefore
invalidates the report.

\subsection{Implemented components and measured result}

The reusable engine is
\texttt{backend/modules/consensus/network\_soak\_qualification.py}.  The
operator entry point is
\texttt{backend/scripts/run\_glyphchain\_network\_soak\_qualification.py}, and
the direct regression is
\texttt{backend/tests/glyphchain/test\_network\_soak\_qualification.py}.  The
retained content-free result is
\texttt{docs/aion/GLYPHCHAIN\_NETWORK\_SOAK\_QUALIFICATION\_REPORT.json}.

The retained run completed all 768 cycles, observed every planned delayed
redelivery, acknowledged every message, treated every repeated put as
idempotent, recovered after every planned restart, validated all partition
journals, returned every network to \texttt{HEALTHY}, rejected the cross-chain
envelope and found no Legal content.  The report state is
\texttt{LOCAL\_BOUNDED\_SOAK\_PASSED} and remains explicitly
\texttt{production\_ready=false}.

Ten direct soak regressions plus the adversarial consensus delay/reordering,
NAT-relay and partition-recovery suites form a focused twenty-five-test set.
All twenty-five pass.

\subsection{Claims and remaining release gates}

The measured duration is diagnostic only.  It is not a transactions-per-second
claim, finality service-level objective or production capacity result.  The
exercise ran on one local computer and did not reproduce carrier NAT, real
internet loss, an independently operated relay, hostile traffic volume or
independent host failure.

Production approval still requires sustained representative load, independent
multi-host restarts, real intermittent networks and relay failures, multiple
independently operated networks, live descriptor-registry/session enforcement,
and independent application-security, consensus-protocol and cryptography
reviews.  Findings from those reviews must be resolved and retested before the
release gate can close.

